% GENERATED FILE. Edit canonical lessons, then run npm run generate:books.
% canonical-source-hash: fnv1a64:1a0f26c4dcdc34de
% canonical-lessons: GE-C08-uhr, GE-C08-uhr-hora, GE-C08-stunde, GE-C08-es-ist-uhr, GE-C08-uhr-ohne-plural, GE-C08-mitte, GE-C08-mittag, GE-C08-mittag-midi, GE-C08-mitternacht, GE-C08-mittag-mitternacht

\chapter{Telling the Time}
\label{ch:time}
\hlchaptermodality{12}

\begin{chapteropening}
\textbf{By the end of this chapter:} \emph{I can ask for and tell the time in German.}

The chapter's thesis is a division of labour. Uhr came from Latin hora with the clocks that needed it, and German already had Stunde for the span, so the loan took the clock-reading job and the native word kept the other. Mittag and Mitternacht are German's own compounds, Mitte with Tag and Nacht, the same thought French assembled out of Latin as midi. It closes on the two frames together: a number with Uhr, or a pivot standing alone.
\end{chapteropening}

\section[die Uhr]{die Uhr — \textquotedblleft{}the clock\textquotedblright{}}

\label{lesson:GE-C08-uhr}



\begin{quote}
Something changes here. The numbers were German's own and the weekdays were German's gods — this word is neither.
\end{quote}

\subsection*{You'll want to know: die Uhr}
\begin{quote}
\textbf{die Uhr} — \textquotedblleft{}the clock,\textquotedblright{} and also \textquotedblleft{}o'clock.\textquotedblright{}
\end{quote}

A \emph{die}-word, and one syllable. It does two jobs: the object on the wall, and the word you say when telling the time.

\begin{sounds}
\begin{itemize}
\raggedright
  \item \texttt{u-long} — a long \emph{u}, the vowel of English \emph{moor}.
  \item \texttt{r-final} — and then a soft \emph{r} that barely closes: German's final \emph{r} colours the vowel and stops rather than rolling.
\end{itemize}

Together: \emph{oor}. One syllable, nothing more.
\end{sounds}

\subsection*{Your turn}
Say these aloud:

\begin{itemize}
\raggedright
  \item \textquotedblleft{}die Uhr\textquotedblright{} — \emph{oor}, long and soft
  \item its two jobs — \textquotedblleft{}the clock; o'clock\textquotedblright{}
  \item it beside a word you own — \textquotedblleft{}der Tag, die Uhr\textquotedblright{}
\end{itemize}

\emph{Twice through:} \textquotedblleft{}die Uhr.\textquotedblright{}

\begin{tcolorbox}[breakable,colback=teal!4,colframe=teal!35!black,title={Before you move on}]
Say \textquotedblleft{}the clock.\textquotedblright{} (\textbf{Die Uhr}.) Which article? (\emph{\textbf{Die}}.) How many syllables? (\textbf{One} — \emph{oor}.) What else does the word mean? (\textbf{O'clock}.) Next: where it came from, which is not where the rest of this book's words come from.
\end{tcolorbox}
\section[Uhr und hōra]{Uhr, taken apart}

\label{lesson:GE-C08-uhr-hora}



\begin{quote}
You have twice met a language keeping some words and buying others. Here is the third time, and it completes a picture.
\end{quote}

\begin{cousinweb}
\emph{Uhr} comes from Latin \emph{\textbf{hōra}}, \textquotedblleft{}hour\textquotedblright{} — the same word that became French \emph{heure}, Italian \emph{ora} and English \textbf{hour}. It reached German through medieval Latin and Low German \emph{ūr(e)}.

Timekeeping arrived with Roman and monastic clocks, and the word rode in with the machinery.
\end{cousinweb}

\begin{culture}
Three layers of the German day, three different origins:

\noindent
\begin{tabularx}{\linewidth}{@{}>{\raggedright\arraybackslash}X>{\raggedright\arraybackslash}X>{\raggedright\arraybackslash}X@{}}
\toprule
\textbf{what} & \textbf{German} & \textbf{where it came from} \\
\midrule
the numbers & \emph{eins, zwei, drei} & native Germanic \\
the weekdays & \emph{Montag, Donnerstag} & Germanic gods \\
the clock & \emph{Uhr} & Latin, from \emph{hōra} \\
\bottomrule
\end{tabularx}

German kept its own counting and its own gods. For the \textbf{machine that measures hours} it reached for Rome.

That is the same rule you were given with the months and again with \emph{Samstag}: the words used every hour of every day are the hardest to displace, and the words that arrive with an institution come in whole. A clock is an institution arriving.
\end{culture}

\subsection*{Your turn}
Say these aloud:

\begin{itemize}
\raggedright
  \item the family — \textquotedblleft{}Uhr, hōra, heure, ora, hour\textquotedblright{}
  \item the three layers — \textquotedblleft{}numbers native, weekdays gods, clock Latin\textquotedblright{}
  \item the reason — \textquotedblleft{}the clock came in, and its word came with it\textquotedblright{}
\end{itemize}

\emph{Twice through:} \textquotedblleft{}Uhr — from hōra.\textquotedblright{}

\begin{tcolorbox}[breakable,colback=teal!4,colframe=teal!35!black,title={Before you move on}]
Which Latin word is \emph{Uhr}? (\emph{\textbf{Hōra}}.) Which English word is its cousin? (\emph{\textbf{Hour}}.) Which two layers of the German day are \textbf{not} borrowed? (\textbf{The numbers and the weekdays}.) Why did this one come in? (\textbf{It arrived with the clocks}.) Next: the German word that was there already.
\end{tcolorbox}
\section[die Stunde]{die Stunde — \textquotedblleft{}an hour\textquotedblright{}}

\label{lesson:GE-C08-stunde}



\begin{quote}
A borrowed word usually arrives because a language needs one. This time German already had a perfectly good word, and took the Latin one anyway.
\end{quote}

\subsection*{You'll want to know: die Stunde}
\begin{quote}
\textbf{die Stunde} — \textquotedblleft{}an hour,\textquotedblright{} in the sense of a \textbf{span of sixty minutes}.
\end{quote}

Native Germanic, and older in the language than \emph{Uhr}. It is the word for a duration: a lesson lasts \emph{eine Stunde}.

\begin{sounds}
\begin{itemize}
\raggedright
  \item \texttt{st-initial-scht} — a German \emph{st-} at the start of a word is \emph{sht-}, not \emph{st-}. So this opens \emph{SHTOON-}, which is the single most useful spelling rule in German after \emph{ei} and \emph{ie}.
  \item \texttt{u-long} — the same long \emph{u} you met in \emph{Uhr}.
\end{itemize}
\end{sounds}

\begin{culture}
The loan did \textbf{not} push the native word out. The two divided the work instead:

\noindent
\begin{tabularx}{\linewidth}{@{}>{\raggedright\arraybackslash}X>{\raggedright\arraybackslash}X@{}}
\toprule
\textbf{for} & \textbf{German uses} \\
\midrule
reading a clock — \textquotedblleft{}o'clock\textquotedblright{} & \emph{Uhr} \\
a length of time — \textquotedblleft{}an hour\textquotedblright{} & \emph{Stunde} \\
\bottomrule
\end{tabularx}

English has the same split and hides it: \emph{hour} covers both, and \emph{o'clock} is a worn-down phrase (\emph{of the clock}) doing the first job.

This is worth holding on to, because it is what borrowing usually looks like. The loan rarely kills the native word; it takes a slice of the meaning and the two settle into different rooms.
\end{culture}

\subsection*{Your turn}
Say these aloud:

\begin{itemize}
\raggedright
  \item \textquotedblleft{}die Stunde\textquotedblright{} — \emph{SHTOON-duh}
  \item the pair and their jobs — \textquotedblleft{}Uhr for the clock, Stunde for the span\textquotedblright{}
  \item the spelling rule — \textquotedblleft{}st- at the front is sht-\textquotedblright{}
\end{itemize}

\emph{Twice through:} \textquotedblleft{}die Uhr, die Stunde.\textquotedblright{}

\begin{tcolorbox}[breakable,colback=teal!4,colframe=teal!35!black,title={Before you move on}]
Say \textquotedblleft{}an hour,\textquotedblright{} meaning a span. (\textbf{Eine Stunde}.) Which word do you use to read a clock? (\emph{\textbf{Uhr}}.) Which of the two is native German? (\emph{\textbf{Stunde}}.) How is \emph{st-} pronounced at the start of a word? (\emph{\textbf{Sht-}}.) Next: putting \emph{Uhr} into a sentence.
\end{tcolorbox}
\section[Es ist zwei Uhr]{Es ist zwei Uhr — telling the time}

\label{lesson:GE-C08-es-ist-uhr}



\begin{quote}
You have the numbers and you have the word. The frame that joins them needs nothing new.
\end{quote}

\begin{grammarlens}[title={Grammar Lens: es ist, a number, Uhr}]
\begin{quote}
\textbf{es ist} + the number + \textbf{Uhr}
\end{quote}

\begin{quote}
\emph{Es ist \textbf{ein} Uhr.} — \textquotedblleft{}It is one o'clock.\textquotedblright{} \emph{Es ist \textbf{zwei} Uhr.} — \textquotedblleft{}It is two o'clock.\textquotedblright{} \emph{Es ist \textbf{zehn} Uhr.} — \textquotedblleft{}It is ten o'clock.\textquotedblright{}
\end{quote}

Word for word that is the English sentence with \emph{o'clock} moved: \emph{it is two o'clock}. Nothing about the order has to be relearned.

One thing to notice: at one o'clock it is \emph{\textbf{ein}} Uhr, not \emph{eins} — the counting form drops its \emph{-s} here exactly as it does before a noun, which you met with \emph{ein} and \emph{eine}.
\end{grammarlens}

\subsection*{Your turn}
Say these aloud:

\begin{itemize}
\raggedright
  \item the frame — \textquotedblleft{}es ist … Uhr\textquotedblright{}
  \item \textquotedblleft{}Es ist zwei Uhr.\textquotedblright{} \textquotedblleft{}Es ist zehn Uhr.\textquotedblright{}
  \item the odd one — \textquotedblleft{}Es ist ein Uhr\textquotedblright{}, not \emph{eins}
  \item the hour it is now, wherever you are
\end{itemize}

\emph{Twice through:} \textquotedblleft{}Es ist zwei Uhr.\textquotedblright{}

\begin{tcolorbox}[breakable,colback=teal!4,colframe=teal!35!black,title={Before you move on}]
Say \textquotedblleft{}it is ten o'clock.\textquotedblright{} (\textbf{Es ist zehn Uhr}.) What are the three pieces? (\emph{\textbf{Es ist}}, \textbf{the number}, \emph{\textbf{Uhr}}.) What happens to \emph{eins} at one o'clock? (\textbf{It drops the \emph{-s}}: \emph{ein Uhr}.) Next: the plural German does not use here.
\end{tcolorbox}
\section[zwei Uhr, nicht zwei Uhren]{zwei Uhr, not zwei Uhren}

\label{lesson:GE-C08-uhr-ohne-plural}



\begin{quote}
Every German noun you have met takes a plural. This one refuses — but only in one of its two jobs.
\end{quote}

\begin{grammarlens}[title={Grammar Lens: the clock-reading Uhr does not count}]
\begin{quote}
Telling the time, \emph{Uhr} stays \textbf{singular} however many hours you name.
\end{quote}

\emph{Zwei Uhr} is \textquotedblleft{}two o'clock.\textquotedblright{} Not \emph{zwei Uhren} — that would be \textbf{two clocks}, the objects on the wall.

\noindent
\begin{tabularx}{\linewidth}{@{}>{\raggedright\arraybackslash}X>{\raggedright\arraybackslash}X@{}}
\toprule
\textbf{German} & \textbf{English} \\
\midrule
\emph{zwei Uhr} & two o'clock \\
\emph{zwei Uhren} & two clocks \\
\bottomrule
\end{tabularx}

English does the identical thing and hides it behind an apostrophe. You would never say \textquotedblleft{}two o'clocks\textquotedblright{} either, because \emph{o'clock} is not a noun being counted — it is a fossilised phrase, \emph{of the clock}. German's \emph{Uhr} is doing that job with no apostrophe to warn you.

So what you are watching is not really about \emph{Uhr}. It is about a word that has stopped being a countable thing and become a \textbf{unit of time}, a job nouns take on in every language and mark differently in each.
\end{grammarlens}

\subsection*{Your turn}
Say these aloud:

\begin{itemize}
\raggedright
  \item the time — \textquotedblleft{}zwei Uhr\textquotedblright{}
  \item the objects — \textquotedblleft{}zwei Uhren\textquotedblright{}
  \item the English that hides it — \textquotedblleft{}two o'clock, two clocks\textquotedblright{}
\end{itemize}

\emph{Twice through:} \textquotedblleft{}zwei Uhr; zwei Uhren.\textquotedblright{}

\begin{tcolorbox}[breakable,colback=teal!4,colframe=teal!35!black,title={Before you move on}]
Say \textquotedblleft{}two o'clock.\textquotedblright{} (\textbf{Zwei Uhr}.) Say \textquotedblleft{}two clocks.\textquotedblright{} (\textbf{Zwei Uhren}.) Which of the two is the plural? (\emph{\textbf{Uhren}}.) Why does the time not take one? (\textbf{It is a unit of time, not a thing being counted}.) Next: the two pivots of the day, and they are German's own.
\end{tcolorbox}
\section[die Mitte]{die Mitte — \textquotedblleft{}the middle\textquotedblright{}}

\label{lesson:GE-C08-mitte}



\begin{quote}
The clock-word was Latin. The two pivots of the German day swing straight back to German's own vocabulary, and they share a piece.
\end{quote}

\subsection*{You'll want to know: die Mitte}
\begin{quote}
\textbf{die Mitte} — \textquotedblleft{}the middle.\textquotedblright{}
\end{quote}

Two syllables, short \emph{i}, and a crisp doubled \emph{t}: \emph{MIT-tuh}.

\begin{cousinweb}
\emph{Mitte} is English \textbf{mid} and \textbf{middle} — the same inherited word, and you can hear it without being told.

Its Latin cousin is \emph{medius}, \textquotedblleft{}middle,\textquotedblright{} which English borrowed several times over: \textbf{medium}, \textbf{median}, \textbf{mediocre} (literally \textquotedblleft{}middling\textquotedblright{}), and the \emph{medi-} of \emph{Mediterranean}, the sea in the middle of the land.

So the everyday Germanic word and the learned Latin one sit side by side in English, and German has the first of them. That is the pattern you have been watching all book: two doors into one root, one worn smooth by daily use and one carried in by scholars.
\end{cousinweb}

\subsection*{Your turn}
Say these aloud:

\begin{itemize}
\raggedright
  \item \textquotedblleft{}die Mitte\textquotedblright{} — \emph{MIT-tuh}
  \item the Germanic side — \textquotedblleft{}Mitte, mid, middle\textquotedblright{}
  \item the Latin side — \textquotedblleft{}medius, medium, median\textquotedblright{}
\end{itemize}

\emph{Twice through:} \textquotedblleft{}die Mitte — the middle.\textquotedblright{}

\begin{tcolorbox}[breakable,colback=teal!4,colframe=teal!35!black,title={Before you move on}]
Say \textquotedblleft{}the middle.\textquotedblright{} (\textbf{Die Mitte}.) Which two English words is it? (\emph{\textbf{Mid}} \textbf{and} \emph{\textbf{middle}}.) Which Latin word is its cousin? (\emph{\textbf{Medius}}.) Name an English word from that side. (\emph{\textbf{Medium}}, \emph{\textbf{median}} or \emph{\textbf{mediocre}}.) Next: the middle of the day.
\end{tcolorbox}
\section[der Mittag]{der Mittag — \textquotedblleft{}noon\textquotedblright{}}

\label{lesson:GE-C08-mittag}



\begin{quote}
You have the middle and you have the day. Put them together before reading on — you will be right.
\end{quote}

\subsection*{You'll want to know: der Mittag}
\begin{quote}
\textbf{der Mittag} — \textquotedblleft{}noon,\textquotedblright{} and literally \textquotedblleft{}\textbf{mid-day}.\textquotedblright{}
\end{quote}

\emph{Mitt-} from \emph{Mitte}, \emph{-tag} from \emph{Tag}. Nothing is hidden and nothing had to be learned: this is the third compound in two chapters you could read on sight, after \emph{Montag} and \emph{Sonnabend}.

English has the identical compound — \textbf{midday} — built from the identical two inherited words. Neither language got it from the other; each still had the pieces and used them.

To say it is noon: \emph{Es ist Mittag.} No number and no \emph{Uhr}, because noon is not an hour you count to.

\subsection*{Your turn}
Say these aloud:

\begin{itemize}
\raggedright
  \item \textquotedblleft{}der Mittag\textquotedblright{}
  \item its two pieces — \textquotedblleft{}Mitte, Tag\textquotedblright{}
  \item the English twin — \textquotedblleft{}Mittag, midday\textquotedblright{}
  \item \textquotedblleft{}Es ist Mittag.\textquotedblright{}
\end{itemize}

\emph{Twice through:} \textquotedblleft{}der Mittag — midday.\textquotedblright{}

\begin{tcolorbox}[breakable,colback=teal!4,colframe=teal!35!black,title={Before you move on}]
Say \textquotedblleft{}noon.\textquotedblright{} (\textbf{Der Mittag}.) What are its two pieces? (\emph{\textbf{Mitte}} \textbf{and} \emph{\textbf{Tag}}.) Say \textquotedblleft{}it is noon.\textquotedblright{} (\textbf{Es ist Mittag}.) Does it need a number? (\textbf{No}.) Next: what French does with the same idea, and why German did not.
\end{tcolorbox}
\section[Mittag und midi]{Mittag and midi}

\label{lesson:GE-C08-mittag-midi}



\begin{quote}
French says \emph{midi} for noon, and it means exactly what \emph{Mittag} means. That is not a borrowing in either direction.
\end{quote}

\begin{cousinweb}
French \emph{midi} comes from Latin \emph{\textbf{medius diēs}} — \textquotedblleft{}middle day.\textquotedblright{} German \emph{Mittag} is \emph{Mitte} plus \emph{Tag} — \textquotedblleft{}middle day.\textquotedblright{}

The \textbf{same thought}, assembled \textbf{twice}, out of different material:

\noindent
\begin{tabularx}{\linewidth}{@{}>{\raggedright\arraybackslash}X>{\raggedright\arraybackslash}X>{\raggedright\arraybackslash}X@{}}
\toprule
\textbf{language} & \textbf{built from} & \textbf{result} \\
\midrule
Latin, then French & \emph{medius} + \emph{diēs} & \emph{midi} \\
German & \emph{Mitte} + \emph{Tag} & \emph{Mittag} \\
\bottomrule
\end{tabularx}

Both families looked at the sun overhead and named the moment for the middle of the day. Neither copied the other; the idea is obvious enough that two languages reached it independently, and each used the words it had.
\end{cousinweb}

\begin{culture}
This is the counterweight to \emph{Uhr}, and the two are worth holding together.

For the \textbf{clock} — a machine that arrived from outside — German took the foreign word whole. For \textbf{noon} — something anyone with a sky already has — German built its own out of pieces it owned.

A language borrows what arrives and builds what was always there. One chapter, both halves of that.
\end{culture}

\subsection*{Your turn}
Say these aloud:

\begin{itemize}
\raggedright
  \item the pair — \textquotedblleft{}Mittag, midi\textquotedblright{}
  \item what each is made of — \textquotedblleft{}Mitte plus Tag; medius plus diēs\textquotedblright{}
  \item the chapter in one line — \textquotedblleft{}borrow the clock, build the noon\textquotedblright{}
\end{itemize}

\emph{Twice through:} \textquotedblleft{}Mittag, midi — the same thought, built twice.\textquotedblright{}

\begin{tcolorbox}[breakable,colback=teal!4,colframe=teal!35!black,title={Before you move on}]
What is French \emph{midi} made of? (\emph{\textbf{Medius}} \textbf{+} \emph{\textbf{diēs}}, \textquotedblleft{}middle day\textquotedblright{}.) And \emph{Mittag}? (\emph{\textbf{Mitte}} \textbf{+} \emph{\textbf{Tag}}.) Did German borrow it? (\textbf{No} — it built its own.) Which word in this chapter \textbf{was} borrowed? (\emph{\textbf{Uhr}}.) Next: the other pivot.
\end{tcolorbox}
\section[die Mitternacht]{die Mitternacht — \textquotedblleft{}midnight\textquotedblright{}}

\label{lesson:GE-C08-mitternacht}



\begin{quote}
The day's other pivot, and the second compound in this chapter you can assemble from words you already have.
\end{quote}

\subsection*{You'll want to know: die Mitternacht}
\begin{quote}
\textbf{die Mitternacht} — \textquotedblleft{}midnight,\textquotedblright{} literally \textquotedblleft{}\textbf{mid-night}.\textquotedblright{}
\end{quote}

\emph{Mitter-} is \emph{Mitte} with the linking form it takes inside a compound, and \emph{-nacht} is the night. English \textbf{midnight} is the same two words again.

To say it: \emph{Es ist Mitternacht.} Like \emph{Mittag}, no number and no \emph{Uhr}.

\begin{sounds}
\begin{itemize}
\raggedright
  \item \texttt{ch-ach} — the raspy back-of-the-throat sound from \emph{Nacht} itself, unchanged inside the longer word.
\end{itemize}

And that \emph{-cht} is the correspondence you were given with the numbers: German \emph{-cht-} answers English \emph{-ght-} every time.

\noindent
\begin{tabularx}{\linewidth}{@{}>{\raggedright\arraybackslash}X>{\raggedright\arraybackslash}X@{}}
\toprule
\textbf{German} & \textbf{English} \\
\midrule
\emph{acht} & eight \\
\emph{Nacht} & night \\
\emph{Licht} & light \\
\bottomrule
\end{tabularx}

You met the first pair counting to ten. Here it is again on a word you now need, which is what a sound correspondence is for — you learn it once and it keeps paying.
\end{sounds}

\subsection*{Your turn}
Say these aloud:

\begin{itemize}
\raggedright
  \item \textquotedblleft{}die Mitternacht\textquotedblright{}
  \item its pieces — \textquotedblleft{}Mitte, Nacht\textquotedblright{}
  \item the correspondence — \textquotedblleft{}acht/eight, Nacht/night, Licht/light\textquotedblright{}
  \item \textquotedblleft{}Es ist Mitternacht.\textquotedblright{}
\end{itemize}

\emph{Twice through:} \textquotedblleft{}Mittag, Mitternacht.\textquotedblright{}

\begin{tcolorbox}[breakable,colback=teal!4,colframe=teal!35!black,title={Before you move on}]
Say \textquotedblleft{}midnight.\textquotedblright{} (\textbf{Die Mitternacht}.) What are its pieces? (\emph{\textbf{Mitte}} \textbf{and} \emph{\textbf{Nacht}}.) What does German \emph{-cht-} answer in English? (\emph{\textbf{-ght-}}.) Give the pair you met with the numbers. (\emph{\textbf{Acht}} \textbf{/} \emph{eight}.) Next: the whole chapter, said through.
\end{tcolorbox}
\section[Practice]{Practice — the clock and the two pivots}

\label{lesson:GE-C08-mittag-mitternacht}



\begin{quote}
No new words. One borrowed word and two built ones, and the difference between them is the chapter. Start by saying the borrowed one, one long syllable — \emph{Uhr} — and then the piece both pivots begin with, \emph{Mitte}.
\end{quote}

\begin{grammarlens}[title={Grammar Lens: three ways to say when}]
\noindent
\begin{tabularx}{\linewidth}{@{}>{\raggedright\arraybackslash}X>{\raggedright\arraybackslash}X@{}}
\toprule
\textbf{when} & \textbf{German} \\
\midrule
one o'clock & \emph{Es ist ein Uhr.} \\
ten o'clock & \emph{Es ist zehn Uhr.} \\
noon & \emph{Es ist Mittag.} \\
midnight & \emph{Es ist Mitternacht.} \\
\bottomrule
\end{tabularx}

One frame, \emph{es ist}, and then either a number with \emph{Uhr} or a pivot on its own. The pivots take no number because they are not hours you count to.
\end{grammarlens}

\begin{culture}
Hold the two halves of the chapter side by side:

\begin{itemize}
\raggedright
  \item \emph{\textbf{Uhr}} came from Latin \emph{hōra}, with the clocks. German had \emph{Stunde} already, so the loan took the clock-reading job and the native word kept the span.
  \item \emph{\textbf{Mittag}} and \emph{\textbf{Mitternacht}} are German's own, built from \emph{Mitte} with \emph{Tag} and \emph{Nacht} — the same thought French assembled out of Latin as \emph{midi}.
\end{itemize}

\textbf{A language borrows what arrives and builds what was always there.} The sky was always there; the clock was delivered.
\end{culture}

\subsection*{Your turn}
Say these aloud:

\begin{itemize}
\raggedright
  \item the four sentences above, in order
  \item the two hour-words and their jobs — \textquotedblleft{}Uhr for the clock, Stunde for the span\textquotedblright{}
  \item the counted and the uncounted — \textquotedblleft{}zwei Uhr; Mittag\textquotedblright{}
  \item the chapter in one line — \textquotedblleft{}borrow the clock, build the noon\textquotedblright{}
\end{itemize}

\emph{Twice through:} \textquotedblleft{}Es ist zwei Uhr. Es ist Mittag. Es ist Mitternacht.\textquotedblright{}

\begin{tcolorbox}[breakable,colback=teal!4,colframe=teal!35!black,title={Before you move on}]
Say \textquotedblleft{}it is two o'clock.\textquotedblright{} (\textbf{Es ist zwei Uhr}.) Say \textquotedblleft{}it is midnight.\textquotedblright{} (\textbf{Es ist Mitternacht}.) Which word here is Latin? (\emph{\textbf{Uhr}}, from \emph{hōra}.) Which two did German build itself? (\emph{\textbf{Mittag}} \textbf{and} \emph{\textbf{Mitternacht}}.) Why does \emph{zwei Uhr} have no plural? (\textbf{It is a unit of time, not a thing being counted}.)
\end{tcolorbox}
